\setauthor{AC}
Im Rahmen der Diplomarbeit wurde mit Mediahub eine Webanwendung konzipiert und umgesetzt, die das Ziel verfolgt, persönliche Streaming-Abonnements zentral zu verwalten und mit Filminformationen zu verknüpfen. Ausgangspunkt war die Beobachtung, dass bei mehreren parallel genutzten Diensten leicht der Überblick über Kosten, Anbieter und tatsächliche Verfügbarkeiten verloren geht. Genau an diesem Problem setzt die Anwendung an, indem sie Verwaltungsfunktionen und filmbezogene Recherche in einer gemeinsamen Lösung zusammenführt.

Im technischen und fachlichen Mittelpunkt meines Projektteils stand dabei die Frontend-Entwicklung. Wesentliche Ergebnisse waren die Umsetzung einer Abo-Verwaltung, die Integration einer Filmsuche mit Detailansicht, die Darstellung passender Anbieterinformationen sowie die Einbindung von Keycloak für geschützte Bereiche. Dadurch konnte eine Benutzeroberfläche entwickelt werden, die nicht nur Informationen anzeigt, sondern einen konkreten praktischen Nutzen für die Zielgruppe erzeugt.

Gleichzeitig zeigte das Projekt, dass die Entwicklung einer solchen Anwendung nicht nur aus der Implementierung einzelner Funktionen besteht. Besonders die Verbindung zwischen Backend-Daten, Benutzeroberfläche, Authentifizierung und verständlicher Benutzerführung erwies sich als anspruchsvoll. Gerade aus diesen Herausforderungen ergaben sich jedoch wichtige Erkenntnisse für die eigene fachliche Weiterentwicklung, insbesondere im Umgang mit moderner Frontend-Architektur, Datenfluss und sicherheitsbezogenen Anforderungen.

Der erreichte Projektstand ist als funktionale und fachlich nachvollziehbare Grundlage zu bewerten. Nicht alle denkbaren Erweiterungen wurden vollständig umgesetzt, was vor allem auf die bewusste Fokussierung auf Filme und auf die Eingrenzung des Projektumfangs zurückzuführen ist. Diese Begrenzung war jedoch sinnvoll, da sie eine klare Zielsetzung und eine realistisch umsetzbare Lösung ermöglichte. Rückblickend zeigte sich zugleich, dass einzelne Entscheidungen in der Frontend-Struktur und bei der frühen Verteilung von Datenlogik suboptimal waren und heute konsequenter entlang klarer fachlicher Verantwortlichkeiten getroffen würden.

Zusammenfassend zeigt die Diplomarbeit, dass eine zentrale Verwaltung von Streaming-Abonnements in Verbindung mit Filminformationen sowohl technisch realisierbar als auch aus Sicht der Nutzenden sinnvoll ist. Mediahub stellt damit eine tragfähige Ausgangsbasis dar, auf der in Zukunft weitere Funktionen und fachliche Erweiterungen aufbauen können. Gleichzeitig bleibt der erreichte Stand bewusst als Zwischenstand zu bewerten: Die Verlaufsfunktion wurde nur teilweise umgesetzt und einzelne Komfortfunktionen der Abo-Verwaltung würden bei einer Weiterentwicklung gezielt ausgebaut werden.
